% Options for packages loaded elsewhere
% Options for packages loaded elsewhere
\PassOptionsToPackage{unicode}{hyperref}
\PassOptionsToPackage{hyphens}{url}
\PassOptionsToPackage{dvipsnames,svgnames,x11names}{xcolor}
%
\documentclass[
  letterpaper,
  DIV=11,
  numbers=noendperiod]{scrartcl}
\usepackage{xcolor}
\usepackage{amsmath,amssymb}
\setcounter{secnumdepth}{5}
\usepackage{iftex}
\ifPDFTeX
  \usepackage[T1]{fontenc}
  \usepackage[utf8]{inputenc}
  \usepackage{textcomp} % provide euro and other symbols
\else % if luatex or xetex
  \usepackage{unicode-math} % this also loads fontspec
  \defaultfontfeatures{Scale=MatchLowercase}
  \defaultfontfeatures[\rmfamily]{Ligatures=TeX,Scale=1}
\fi
\usepackage{lmodern}
\ifPDFTeX\else
  % xetex/luatex font selection
\fi
% Use upquote if available, for straight quotes in verbatim environments
\IfFileExists{upquote.sty}{\usepackage{upquote}}{}
\IfFileExists{microtype.sty}{% use microtype if available
  \usepackage[]{microtype}
  \UseMicrotypeSet[protrusion]{basicmath} % disable protrusion for tt fonts
}{}
\makeatletter
\@ifundefined{KOMAClassName}{% if non-KOMA class
  \IfFileExists{parskip.sty}{%
    \usepackage{parskip}
  }{% else
    \setlength{\parindent}{0pt}
    \setlength{\parskip}{6pt plus 2pt minus 1pt}}
}{% if KOMA class
  \KOMAoptions{parskip=half}}
\makeatother
% Make \paragraph and \subparagraph free-standing
\makeatletter
\ifx\paragraph\undefined\else
  \let\oldparagraph\paragraph
  \renewcommand{\paragraph}{
    \@ifstar
      \xxxParagraphStar
      \xxxParagraphNoStar
  }
  \newcommand{\xxxParagraphStar}[1]{\oldparagraph*{#1}\mbox{}}
  \newcommand{\xxxParagraphNoStar}[1]{\oldparagraph{#1}\mbox{}}
\fi
\ifx\subparagraph\undefined\else
  \let\oldsubparagraph\subparagraph
  \renewcommand{\subparagraph}{
    \@ifstar
      \xxxSubParagraphStar
      \xxxSubParagraphNoStar
  }
  \newcommand{\xxxSubParagraphStar}[1]{\oldsubparagraph*{#1}\mbox{}}
  \newcommand{\xxxSubParagraphNoStar}[1]{\oldsubparagraph{#1}\mbox{}}
\fi
\makeatother


\usepackage{longtable,booktabs,array}
\usepackage{calc} % for calculating minipage widths
% Correct order of tables after \paragraph or \subparagraph
\usepackage{etoolbox}
\makeatletter
\patchcmd\longtable{\par}{\if@noskipsec\mbox{}\fi\par}{}{}
\makeatother
% Allow footnotes in longtable head/foot
\IfFileExists{footnotehyper.sty}{\usepackage{footnotehyper}}{\usepackage{footnote}}
\makesavenoteenv{longtable}
\usepackage{graphicx}
\makeatletter
\newsavebox\pandoc@box
\newcommand*\pandocbounded[1]{% scales image to fit in text height/width
  \sbox\pandoc@box{#1}%
  \Gscale@div\@tempa{\textheight}{\dimexpr\ht\pandoc@box+\dp\pandoc@box\relax}%
  \Gscale@div\@tempb{\linewidth}{\wd\pandoc@box}%
  \ifdim\@tempb\p@<\@tempa\p@\let\@tempa\@tempb\fi% select the smaller of both
  \ifdim\@tempa\p@<\p@\scalebox{\@tempa}{\usebox\pandoc@box}%
  \else\usebox{\pandoc@box}%
  \fi%
}
% Set default figure placement to htbp
\def\fps@figure{htbp}
\makeatother





\setlength{\emergencystretch}{3em} % prevent overfull lines

\providecommand{\tightlist}{%
  \setlength{\itemsep}{0pt}\setlength{\parskip}{0pt}}



 


\KOMAoption{captions}{tableheading}
\makeatletter
\@ifpackageloaded{caption}{}{\usepackage{caption}}
\AtBeginDocument{%
\ifdefined\contentsname
  \renewcommand*\contentsname{Table of contents}
\else
  \newcommand\contentsname{Table of contents}
\fi
\ifdefined\listfigurename
  \renewcommand*\listfigurename{List of Figures}
\else
  \newcommand\listfigurename{List of Figures}
\fi
\ifdefined\listtablename
  \renewcommand*\listtablename{List of Tables}
\else
  \newcommand\listtablename{List of Tables}
\fi
\ifdefined\figurename
  \renewcommand*\figurename{Figure}
\else
  \newcommand\figurename{Figure}
\fi
\ifdefined\tablename
  \renewcommand*\tablename{Table}
\else
  \newcommand\tablename{Table}
\fi
}
\@ifpackageloaded{float}{}{\usepackage{float}}
\floatstyle{ruled}
\@ifundefined{c@chapter}{\newfloat{codelisting}{h}{lop}}{\newfloat{codelisting}{h}{lop}[chapter]}
\floatname{codelisting}{Listing}
\newcommand*\listoflistings{\listof{codelisting}{List of Listings}}
\makeatother
\makeatletter
\makeatother
\makeatletter
\@ifpackageloaded{caption}{}{\usepackage{caption}}
\@ifpackageloaded{subcaption}{}{\usepackage{subcaption}}
\makeatother
\usepackage{bookmark}
\IfFileExists{xurl.sty}{\usepackage{xurl}}{} % add URL line breaks if available
\urlstyle{same}
\hypersetup{
  pdftitle={DORA},
  pdfauthor={Joao Luiz D'Andrea},
  colorlinks=true,
  linkcolor={blue},
  filecolor={Maroon},
  citecolor={Blue},
  urlcolor={Blue},
  pdfcreator={LaTeX via pandoc}}


\title{DORA}
\usepackage{etoolbox}
\makeatletter
\providecommand{\subtitle}[1]{% add subtitle to \maketitle
  \apptocmd{\@title}{\par {\large #1 \par}}{}{}
}
\makeatother
\subtitle{Digital Operational Resilience Act --- Governance, Risk,
Controls, and Implementation Template}
\author{Joao Luiz D'Andrea}
\date{}
\begin{document}
\maketitle

\renewcommand*\contentsname{Table of contents}
{
\hypersetup{linkcolor=}
\setcounter{tocdepth}{3}
\tableofcontents
}

\section{Executive Summary}\label{executive-summary}

The Digital Operational Resilience Act (DORA) establishes a harmonised
European Union framework for managing digital operational resilience
across the financial sector. It strengthens the ability of financial
entities to withstand, respond to, and recover from ICT-related
disruptions, cyber incidents, and third-party service failures.

This document provides a practical template for defining the
organisation's DORA governance model, ICT risk management approach,
control structure, implementation roadmap, and assurance model. It can
be used to document how the organisation aligns resilience,
cybersecurity, third-party risk, and regulatory compliance into a
coherent operating model.

\section{Overview}\label{overview}

\subsection{Purpose}\label{purpose}

This document provides a practical template for defining, implementing,
operating, and assuring compliance with DORA across the enterprise.

\subsection{Objectives}\label{objectives}

The main objectives of applying DORA are to:

\begin{itemize}
\tightlist
\item
  strengthen digital operational resilience across regulated financial
  services
\item
  establish clear accountability for ICT risk and resilience
\item
  improve prevention, detection, response, and recovery capabilities
\item
  ensure consistent handling and reporting of ICT-related incidents
\item
  strengthen oversight of ICT third-party dependencies and concentration
  risk
\item
  support resilience testing and continuous improvement
\item
  integrate operational resilience with enterprise governance, risk,
  compliance, and assurance
\end{itemize}

\subsection{Scope}\label{scope}

Describe the scope of DORA applicability in the organisation.

\subsubsection{In scope}\label{in-scope}

\begin{itemize}
\tightlist
\item
  regulated financial entities and relevant business functions
\item
  ICT services supporting critical or important functions
\item
  applications, infrastructure, cloud services, and networks
\item
  cyber resilience, ICT operations, and service continuity arrangements
\item
  ICT third-party service providers and contractual arrangements
\item
  digital operational resilience testing activities
\item
  incident management and regulatory reporting processes
\end{itemize}

\subsubsection{Out of scope}\label{out-of-scope}

Document justified exclusions, if any.

\subsection{Intended Audience}\label{intended-audience}

\begin{itemize}
\tightlist
\item
  Board and executive leadership
\item
  Chief Information Officer
\item
  Chief Information Security Officer
\item
  Chief Risk Officer
\item
  Operational resilience teams
\item
  Technology and operations teams
\item
  Procurement and vendor management teams
\item
  Legal and compliance teams
\item
  Internal audit and assurance teams
\item
  PMO / transformation teams
\end{itemize}

\section{Regulatory Foundation}\label{regulatory-foundation}

\subsection{Background}\label{background}

DORA provides a unified regulatory framework for digital operational
resilience in the financial sector. It addresses the management of ICT
risk, resilience testing, ICT incident handling and reporting, ICT
third-party risk, and supervisory oversight of critical ICT third-party
providers.

\subsection{Regulatory Objectives}\label{regulatory-objectives}

Document how the organisation interprets the regulatory objectives of
DORA:

\begin{itemize}
\tightlist
\item
  maintain resilience of ICT systems supporting business processes
\item
  reduce the impact of ICT-related disruptions and cyber threats
\item
  strengthen governance and accountability for ICT risk
\item
  improve visibility and management of third-party dependencies
\item
  increase consistency in resilience testing and incident reporting
\item
  enhance financial sector-wide resilience
\end{itemize}

\subsection{Core DORA Pillars}\label{core-dora-pillars}

Describe how the organisation addresses the main DORA domains:

\begin{itemize}
\tightlist
\item
  ICT risk management
\item
  ICT-related incident management, classification, and reporting
\item
  digital operational resilience testing
\item
  ICT third-party risk management
\item
  information-sharing arrangements
\end{itemize}

\subsection{Applicability and Entity
Context}\label{applicability-and-entity-context}

Document the specific applicability of DORA to the organisation,
including:

\begin{itemize}
\tightlist
\item
  legal entity structure
\item
  regulatory status
\item
  business model
\item
  group / subsidiary context
\item
  jurisdictional footprint
\item
  critical and important functions
\item
  supervisory relationships
\end{itemize}

\section{Governance Model}\label{governance-model}

\subsection{Governance Objectives}\label{governance-objectives}

Define the governance objectives for DORA compliance:

\begin{itemize}
\tightlist
\item
  establish board-level oversight of digital operational resilience
\item
  define accountability for ICT risk and resilience decisions
\item
  align resilience priorities with enterprise risk appetite and business
  continuity objectives
\item
  embed DORA requirements into operations, architecture, procurement,
  and change
\item
  maintain evidence of compliance and control effectiveness
\item
  support continuous improvement through testing, monitoring, assurance,
  and remediation
\end{itemize}

\subsection{Governance Structure}\label{governance-structure}

Describe the governance bodies and their responsibilities.

\subsubsection{Example governance
bodies}\label{example-governance-bodies}

\begin{itemize}
\tightlist
\item
  Board / Risk Committee
\item
  Executive Operational Resilience Committee
\item
  Technology and Security Steering Committee
\item
  ICT Risk Committee
\item
  Third-Party Risk Committee
\item
  Architecture Review Board
\item
  Incident Review Board
\item
  Internal Audit
\end{itemize}

\subsection{Three Lines Model}\label{three-lines-model}

Describe responsibilities across the three lines.

\subsubsection{First line}\label{first-line}

\begin{itemize}
\tightlist
\item
  ownership of business services, ICT operations, and controls
\item
  implementation of resilience, security, and continuity controls
\item
  incident handling and service restoration
\item
  evidence production and remediation execution
\end{itemize}

\subsubsection{Second line}\label{second-line}

\begin{itemize}
\tightlist
\item
  ICT risk oversight
\item
  resilience policy and standards
\item
  compliance monitoring
\item
  independent challenge
\item
  regulatory coordination
\end{itemize}

\subsubsection{Third line}\label{third-line}

\begin{itemize}
\tightlist
\item
  internal audit
\item
  independent assurance
\item
  control effectiveness review
\end{itemize}

\section{Roles and Responsibilities}\label{roles-and-responsibilities}

\subsection{Key Roles}\label{key-roles}

Document the main DORA-related roles.

\begin{itemize}
\tightlist
\item
  Board
\item
  CEO / Executive Management
\item
  Chief Information Officer
\item
  Chief Information Security Officer
\item
  Chief Risk Officer
\item
  Chief Operating Officer
\item
  Operational Resilience Lead
\item
  ICT Risk Manager
\item
  Incident Response Lead
\item
  Service Owner
\item
  Infrastructure / Platform Owner
\item
  Third-Party Risk Manager
\item
  Procurement Lead
\item
  Legal and Compliance Lead
\item
  Internal Audit
\end{itemize}

\subsection{RACI Matrix}\label{raci-matrix}

Use the following model and adapt it.

\begin{longtable}[]{@{}
  >{\raggedright\arraybackslash}p{(\linewidth - 18\tabcolsep) * \real{0.1000}}
  >{\raggedright\arraybackslash}p{(\linewidth - 18\tabcolsep) * \real{0.1000}}
  >{\raggedright\arraybackslash}p{(\linewidth - 18\tabcolsep) * \real{0.1000}}
  >{\raggedright\arraybackslash}p{(\linewidth - 18\tabcolsep) * \real{0.1000}}
  >{\raggedright\arraybackslash}p{(\linewidth - 18\tabcolsep) * \real{0.1000}}
  >{\raggedright\arraybackslash}p{(\linewidth - 18\tabcolsep) * \real{0.1000}}
  >{\raggedright\arraybackslash}p{(\linewidth - 18\tabcolsep) * \real{0.1000}}
  >{\raggedright\arraybackslash}p{(\linewidth - 18\tabcolsep) * \real{0.1000}}
  >{\raggedright\arraybackslash}p{(\linewidth - 18\tabcolsep) * \real{0.1000}}
  >{\raggedright\arraybackslash}p{(\linewidth - 18\tabcolsep) * \real{0.1000}}@{}}
\toprule\noalign{}
\begin{minipage}[b]{\linewidth}\raggedright
Activity
\end{minipage} & \begin{minipage}[b]{\linewidth}\raggedright
Board
\end{minipage} & \begin{minipage}[b]{\linewidth}\raggedright
Exec
\end{minipage} & \begin{minipage}[b]{\linewidth}\raggedright
CIO
\end{minipage} & \begin{minipage}[b]{\linewidth}\raggedright
CISO
\end{minipage} & \begin{minipage}[b]{\linewidth}\raggedright
CRO
\end{minipage} & \begin{minipage}[b]{\linewidth}\raggedright
Ops Resilience
\end{minipage} & \begin{minipage}[b]{\linewidth}\raggedright
Service Owner
\end{minipage} & \begin{minipage}[b]{\linewidth}\raggedright
Third-Party Risk
\end{minipage} & \begin{minipage}[b]{\linewidth}\raggedright
Audit
\end{minipage} \\
\midrule\noalign{}
\endhead
\bottomrule\noalign{}
\endlastfoot
Approve digital operational resilience strategy & A & R & C & C & C & C
& I & I & I \\
Maintain ICT risk framework & I & A & R & R & C & C & C & I & I \\
Classify and report major ICT incidents & I & A & R & R & C & C & C & I
& I \\
Maintain testing programme & I & A & C & C & C & R & C & I & I \\
Manage ICT third-party risk & I & A & C & C & C & C & C & R & I \\
Maintain register of information & I & A & C & C & C & C & I & R & I \\
Oversee remediation actions & I & A & R & R & C & C & R & C & I \\
Perform independent assurance & I & I & I & I & I & I & I & I & R/A \\
\end{longtable}

\section{Policy Framework}\label{policy-framework}

\subsection{Core Policies}\label{core-policies}

List and describe the main DORA-related policies.

\begin{itemize}
\tightlist
\item
  Digital Operational Resilience Policy
\item
  ICT Risk Management Policy
\item
  Information Security Policy
\item
  ICT Incident Management Policy
\item
  Business Continuity and Disaster Recovery Policy
\item
  Operational Resilience Testing Policy
\item
  Third-Party Risk Management Policy
\item
  Outsourcing / ICT Provider Management Policy
\item
  Change and Release Management Policy
\item
  Asset and Configuration Management Policy
\item
  Logging and Monitoring Standard
\item
  Crisis Management and Communications Policy
\end{itemize}

\subsection{Standards and Procedures}\label{standards-and-procedures}

Document the standards, procedures, playbooks, and operational guidance
supporting compliance.

\section{Core Processes}\label{core-processes}

\subsection{ICT Risk Management}\label{ict-risk-management}

Describe the process for managing ICT risk across the lifecycle:

\begin{itemize}
\tightlist
\item
  governance and accountability
\item
  ICT asset identification
\item
  risk assessment
\item
  protection and prevention
\item
  detection
\item
  response
\item
  recovery
\item
  lessons learned
\item
  continuous improvement
\end{itemize}

\subsection{ICT-Related Incident
Management}\label{ict-related-incident-management}

Describe the incident lifecycle:

\begin{itemize}
\tightlist
\item
  event detection
\item
  triage
\item
  classification
\item
  escalation
\item
  containment
\item
  eradication
\item
  recovery
\item
  root cause analysis
\item
  lessons learned
\end{itemize}

\subsection{Incident Classification and
Reporting}\label{incident-classification-and-reporting}

Document how the organisation handles:

\begin{itemize}
\tightlist
\item
  incident classification criteria
\item
  major incident thresholds
\item
  internal escalation
\item
  competent authority notifications
\item
  intermediate reporting
\item
  final reporting
\item
  post-incident review
\item
  evidence retention
\end{itemize}

\subsection{Digital Operational Resilience
Testing}\label{digital-operational-resilience-testing}

Describe the testing model for:

\begin{itemize}
\tightlist
\item
  annual testing planning
\item
  scenario-based testing
\item
  vulnerability assessments
\item
  penetration testing
\item
  business continuity testing
\item
  disaster recovery exercises
\item
  crisis management exercises
\item
  threat-led penetration testing where applicable
\item
  remediation tracking
\end{itemize}

\subsection{ICT Third-Party Risk
Management}\label{ict-third-party-risk-management}

Describe controls for:

\begin{itemize}
\tightlist
\item
  due diligence
\item
  onboarding assessments
\item
  contractual requirements
\item
  criticality assessment
\item
  concentration risk assessment
\item
  ongoing monitoring
\item
  exit planning
\item
  substitution strategy
\item
  assurance reviews
\end{itemize}

\subsection{Register of Information}\label{register-of-information}

Document how the organisation maintains a complete and accurate register
of information for contractual arrangements with ICT third-party service
providers.

\subsubsection{Minimum content}\label{minimum-content}

\begin{itemize}
\tightlist
\item
  supplier identity
\item
  service description
\item
  supported business entity
\item
  supported function
\item
  critical or important function flag
\item
  location and jurisdiction
\item
  subcontracting arrangements
\item
  contract dates
\item
  renewal terms
\item
  exit clauses
\item
  resilience and security obligations
\end{itemize}

\subsection{Information Sharing
Arrangements}\label{information-sharing-arrangements}

Describe how the organisation governs participation in cyber threat
intelligence or sectoral information-sharing arrangements, including:

\begin{itemize}
\tightlist
\item
  legal and compliance approval
\item
  confidentiality controls
\item
  participation criteria
\item
  information handling requirements
\item
  oversight and review
\end{itemize}

\section{Control Framework}\label{control-framework}

\subsection{Control Domains}\label{control-domains}

Organise the DORA control environment into domains.

\subsubsection{Example control domains}\label{example-control-domains}

\begin{itemize}
\tightlist
\item
  governance and oversight
\item
  ICT risk management
\item
  asset and dependency management
\item
  protection and prevention
\item
  monitoring and detection
\item
  ICT incident response
\item
  resilience and recovery
\item
  testing and validation
\item
  ICT third-party risk
\item
  concentration risk management
\item
  regulatory reporting
\item
  evidence and assurance
\end{itemize}

\subsection{Sample Control Catalogue}\label{sample-control-catalogue}

\begin{longtable}[]{@{}
  >{\raggedright\arraybackslash}p{(\linewidth - 10\tabcolsep) * \real{0.1667}}
  >{\raggedright\arraybackslash}p{(\linewidth - 10\tabcolsep) * \real{0.1667}}
  >{\raggedright\arraybackslash}p{(\linewidth - 10\tabcolsep) * \real{0.1667}}
  >{\raggedright\arraybackslash}p{(\linewidth - 10\tabcolsep) * \real{0.1667}}
  >{\raggedright\arraybackslash}p{(\linewidth - 10\tabcolsep) * \real{0.1667}}
  >{\raggedright\arraybackslash}p{(\linewidth - 10\tabcolsep) * \real{0.1667}}@{}}
\toprule\noalign{}
\begin{minipage}[b]{\linewidth}\raggedright
Control ID
\end{minipage} & \begin{minipage}[b]{\linewidth}\raggedright
Control Name
\end{minipage} & \begin{minipage}[b]{\linewidth}\raggedright
Objective
\end{minipage} & \begin{minipage}[b]{\linewidth}\raggedright
Control Owner
\end{minipage} & \begin{minipage}[b]{\linewidth}\raggedright
Frequency
\end{minipage} & \begin{minipage}[b]{\linewidth}\raggedright
Evidence
\end{minipage} \\
\midrule\noalign{}
\endhead
\bottomrule\noalign{}
\endlastfoot
DORA-01 & ICT risk framework approved & Ensure governance and
accountability & CIO / CRO & Annual & Approved framework \\
DORA-02 & Major incident process maintained & Ensure timely incident
handling and escalation & CISO & Annual review & Incident procedure \\
DORA-03 & Testing programme executed & Validate resilience capabilities
& Resilience Lead & Annual / quarterly & Test reports \\
DORA-04 & ICT third-party due diligence performed & Reduce supplier risk
exposure & Third-Party Risk Manager & Onboarding / annual & Due
diligence records \\
DORA-05 & Register of information maintained & Ensure complete supplier
transparency & Vendor Management Lead & Ongoing & Register of
information \\
DORA-06 & Recovery exercises completed & Validate restoration capability
& Operations / Resilience Lead & Annual & Exercise evidence \\
\end{longtable}

\section{Artefacts and Evidence}\label{artefacts-and-evidence}

\subsection{Required Artefacts}\label{required-artefacts}

Document the key DORA artefacts.

\begin{itemize}
\tightlist
\item
  digital operational resilience strategy
\item
  ICT risk management framework
\item
  ICT asset inventory
\item
  service dependency maps
\item
  critical and important functions register
\item
  ICT risk register
\item
  incident management procedure
\item
  major incident reporting procedure
\item
  incident register
\item
  testing policy and testing plan
\item
  test scripts and exercise reports
\item
  remediation tracker
\item
  ICT third-party register
\item
  register of information
\item
  due diligence assessments
\item
  contractual clause library
\item
  exit plans
\item
  business continuity plans
\item
  disaster recovery plans
\item
  crisis communication plans
\item
  management information and dashboard packs
\item
  audit reports
\end{itemize}

\subsection{Evidence Model}\label{evidence-model}

For each artefact define:

\begin{itemize}
\tightlist
\item
  owner
\item
  purpose
\item
  location
\item
  frequency of update
\item
  retention period
\item
  assurance mechanism
\end{itemize}

\section{Architecture and Technology
Implications}\label{architecture-and-technology-implications}

\subsection{Resilience by Design}\label{resilience-by-design}

Describe how DORA requirements are embedded into:

\begin{itemize}
\tightlist
\item
  enterprise architecture
\item
  solution design
\item
  cloud adoption
\item
  infrastructure engineering
\item
  network design
\item
  identity and access management
\item
  observability and monitoring
\item
  backup and recovery architecture
\item
  high availability and failover design
\item
  supplier integration patterns
\end{itemize}

\subsection{Technology Landscape View}\label{technology-landscape-view}

Document the resilience-relevant technology landscape:

\begin{itemize}
\tightlist
\item
  critical business services
\item
  applications and platforms
\item
  infrastructure and cloud services
\item
  networks and communications
\item
  data stores and backup services
\item
  cyber tooling
\item
  service desk and operations tooling
\item
  third-party technology dependencies
\end{itemize}

\subsection{Dependency Mapping}\label{dependency-mapping}

Include or reference:

\begin{itemize}
\tightlist
\item
  business service maps
\item
  application dependency diagrams
\item
  infrastructure dependency diagrams
\item
  data flow diagrams
\item
  third-party dependency maps
\item
  concentration risk views
\item
  recovery dependency maps
\end{itemize}

\section{Risk Management}\label{risk-management}

\subsection{DORA Risk Model}\label{dora-risk-model}

Describe how ICT and resilience risks are identified, assessed, treated,
monitored, and escalated.

\subsection{Example Risk Categories}\label{example-risk-categories}

\begin{itemize}
\tightlist
\item
  weak board oversight of ICT risk
\item
  incomplete ICT asset visibility
\item
  inadequate resilience controls
\item
  insufficient monitoring and detection
\item
  delayed incident escalation
\item
  poor recovery capability
\item
  excessive third-party concentration
\item
  weak contractual control over suppliers
\item
  inaccurate register of information
\item
  lack of evidence of testing and remediation
\end{itemize}

\subsection{Risk Register Template}\label{risk-register-template}

\begin{longtable}[]{@{}
  >{\raggedright\arraybackslash}p{(\linewidth - 16\tabcolsep) * \real{0.1111}}
  >{\raggedright\arraybackslash}p{(\linewidth - 16\tabcolsep) * \real{0.1111}}
  >{\raggedright\arraybackslash}p{(\linewidth - 16\tabcolsep) * \real{0.1111}}
  >{\raggedright\arraybackslash}p{(\linewidth - 16\tabcolsep) * \real{0.1111}}
  >{\raggedright\arraybackslash}p{(\linewidth - 16\tabcolsep) * \real{0.1111}}
  >{\raggedright\arraybackslash}p{(\linewidth - 16\tabcolsep) * \real{0.1111}}
  >{\raggedright\arraybackslash}p{(\linewidth - 16\tabcolsep) * \real{0.1111}}
  >{\raggedright\arraybackslash}p{(\linewidth - 16\tabcolsep) * \real{0.1111}}
  >{\raggedright\arraybackslash}p{(\linewidth - 16\tabcolsep) * \real{0.1111}}@{}}
\toprule\noalign{}
\begin{minipage}[b]{\linewidth}\raggedright
Risk ID
\end{minipage} & \begin{minipage}[b]{\linewidth}\raggedright
Risk Description
\end{minipage} & \begin{minipage}[b]{\linewidth}\raggedright
Cause
\end{minipage} & \begin{minipage}[b]{\linewidth}\raggedright
Impact
\end{minipage} & \begin{minipage}[b]{\linewidth}\raggedright
Likelihood
\end{minipage} & \begin{minipage}[b]{\linewidth}\raggedright
Inherent Risk
\end{minipage} & \begin{minipage}[b]{\linewidth}\raggedright
Control Strength
\end{minipage} & \begin{minipage}[b]{\linewidth}\raggedright
Residual Risk
\end{minipage} & \begin{minipage}[b]{\linewidth}\raggedright
Owner
\end{minipage} \\
\midrule\noalign{}
\endhead
\bottomrule\noalign{}
\endlastfoot
DR-01 & Example & Example & Example & Medium & High & Medium & Medium &
Example \\
\end{longtable}

\section{Metrics and Reporting}\label{metrics-and-reporting}

\subsection{Key Metrics}\label{key-metrics}

Examples of DORA metrics:

\begin{itemize}
\tightlist
\item
  percentage of critical services with documented dependencies
\item
  percentage of ICT assets covered by resilience controls
\item
  number of major ICT incidents
\item
  mean time to detect
\item
  mean time to contain
\item
  mean time to recover
\item
  percentage of planned resilience tests completed
\item
  percentage of severe findings remediated on time
\item
  percentage of critical suppliers assessed
\item
  concentration risk exposure by provider
\item
  percentage completeness of register of information
\item
  number of open audit or regulatory findings
\end{itemize}

\subsection{Dashboard Structure}\label{dashboard-structure}

Suggested dashboard views:

\begin{itemize}
\tightlist
\item
  executive resilience dashboard
\item
  ICT risk dashboard
\item
  incident and recovery dashboard
\item
  testing and remediation dashboard
\item
  third-party and concentration risk dashboard
\item
  audit and compliance dashboard
\end{itemize}

\section{Audit and Assurance}\label{audit-and-assurance}

\subsection{Assurance Model}\label{assurance-model}

Document the assurance approach across:

\begin{itemize}
\tightlist
\item
  management self-assessment
\item
  second-line compliance reviews
\item
  control testing
\item
  scenario and exercise validation
\item
  internal audit
\item
  supervisory review support
\end{itemize}

\subsection{Testing Themes}\label{testing-themes}

\begin{itemize}
\tightlist
\item
  governance and accountability
\item
  ICT risk framework effectiveness
\item
  incident classification and escalation
\item
  reporting timeliness and quality
\item
  resilience testing coverage
\item
  remediation effectiveness
\item
  third-party governance
\item
  register of information quality
\item
  recovery readiness
\item
  evidence completeness
\end{itemize}

\section{Implementation Roadmap}\label{implementation-roadmap}

\subsection{Maturity Phases}\label{maturity-phases}

\subsubsection{Phase 1: Foundation}\label{phase-1-foundation}

\begin{itemize}
\tightlist
\item
  establish governance and accountability
\item
  define DORA policy framework
\item
  identify critical and important functions
\item
  establish asset and dependency baseline
\item
  define incident classification and reporting model
\end{itemize}

\subsubsection{Phase 2: Control Design}\label{phase-2-control-design}

\begin{itemize}
\tightlist
\item
  strengthen ICT risk management controls
\item
  implement resilience testing framework
\item
  formalise third-party governance and register of information
\item
  improve monitoring, response, and recovery controls
\end{itemize}

\subsubsection{Phase 3: Integration}\label{phase-3-integration}

\begin{itemize}
\tightlist
\item
  integrate DORA with enterprise risk, architecture, and procurement
\item
  improve dependency mapping and concentration risk management
\item
  align service continuity, cyber resilience, and crisis management
\item
  automate evidence collection and reporting workflows
\end{itemize}

\subsubsection{Phase 4: Optimisation}\label{phase-4-optimisation}

\begin{itemize}
\tightlist
\item
  mature dashboards and management information
\item
  strengthen scenario-based and threat-led exercises
\item
  improve remediation governance
\item
  enhance continuous monitoring and assurance
\end{itemize}

\subsection{Roadmap Table}\label{roadmap-table}

\begin{longtable}[]{@{}
  >{\raggedright\arraybackslash}p{(\linewidth - 12\tabcolsep) * \real{0.1429}}
  >{\raggedright\arraybackslash}p{(\linewidth - 12\tabcolsep) * \real{0.1429}}
  >{\raggedright\arraybackslash}p{(\linewidth - 12\tabcolsep) * \real{0.1429}}
  >{\raggedright\arraybackslash}p{(\linewidth - 12\tabcolsep) * \real{0.1429}}
  >{\raggedright\arraybackslash}p{(\linewidth - 12\tabcolsep) * \real{0.1429}}
  >{\raggedright\arraybackslash}p{(\linewidth - 12\tabcolsep) * \real{0.1429}}
  >{\raggedright\arraybackslash}p{(\linewidth - 12\tabcolsep) * \real{0.1429}}@{}}
\toprule\noalign{}
\begin{minipage}[b]{\linewidth}\raggedright
Initiative
\end{minipage} & \begin{minipage}[b]{\linewidth}\raggedright
Description
\end{minipage} & \begin{minipage}[b]{\linewidth}\raggedright
Owner
\end{minipage} & \begin{minipage}[b]{\linewidth}\raggedright
Priority
\end{minipage} & \begin{minipage}[b]{\linewidth}\raggedright
Timeline
\end{minipage} & \begin{minipage}[b]{\linewidth}\raggedright
Dependency
\end{minipage} & \begin{minipage}[b]{\linewidth}\raggedright
Status
\end{minipage} \\
\midrule\noalign{}
\endhead
\bottomrule\noalign{}
\endlastfoot
Example & Example & Example & High & Q1-Q2 & Example & Planned \\
\end{longtable}

\section{Cross-Framework Integration}\label{cross-framework-integration}

\subsection{Relationship to Other
Frameworks}\label{relationship-to-other-frameworks}

Describe how DORA integrates with:

\begin{itemize}
\tightlist
\item
  NIS2
\item
  ISO/IEC 27001
\item
  GDPR
\item
  AI Act
\item
  NIST Cybersecurity Framework
\item
  COBIT
\item
  ITIL
\item
  TOGAF
\item
  enterprise risk management
\item
  business continuity and operational resilience frameworks
\end{itemize}

\subsection{Example Mapping Themes}\label{example-mapping-themes}

\begin{longtable}[]{@{}
  >{\raggedright\arraybackslash}p{(\linewidth - 4\tabcolsep) * \real{0.3333}}
  >{\raggedright\arraybackslash}p{(\linewidth - 4\tabcolsep) * \real{0.3333}}
  >{\raggedright\arraybackslash}p{(\linewidth - 4\tabcolsep) * \real{0.3333}}@{}}
\toprule\noalign{}
\begin{minipage}[b]{\linewidth}\raggedright
DORA Theme
\end{minipage} & \begin{minipage}[b]{\linewidth}\raggedright
Related Framework
\end{minipage} & \begin{minipage}[b]{\linewidth}\raggedright
Integration Point
\end{minipage} \\
\midrule\noalign{}
\endhead
\bottomrule\noalign{}
\endlastfoot
ICT risk management & ISO 27001 / NIST CSF & Cybersecurity and control
framework \\
Incident reporting & NIS2 / GDPR & Regulatory notification and incident
governance \\
Third-party risk & COBIT / ISO 27001 & Supplier and outsourcing
governance \\
Resilience testing & NIST CSF / BCM & Validation of resilience
capabilities \\
Architecture resilience & TOGAF & Design authority and resilience by
design \\
\end{longtable}

\section{Maturity Model}\label{maturity-model}

\subsection{Suggested Maturity Levels}\label{suggested-maturity-levels}

\subsubsection{Level 1: Initial}\label{level-1-initial}

Resilience practices are reactive, fragmented, and poorly evidenced.

\subsubsection{Level 2: Repeatable}\label{level-2-repeatable}

Basic ICT risk, continuity, and incident processes exist but are
inconsistently applied.

\subsubsection{Level 3: Defined}\label{level-3-defined}

Policies, roles, testing, and third-party controls are documented and
implemented across key areas.

\subsubsection{Level 4: Managed}\label{level-4-managed}

Resilience capabilities are measured, risk-based, tested, and governed
through regular reporting.

\subsubsection{Level 5: Optimised}\label{level-5-optimised}

Digital operational resilience is embedded across governance,
architecture, operations, and supply chain management with strong
assurance and continuous improvement.

\subsection{Assessment Dimensions}\label{assessment-dimensions}

\begin{itemize}
\tightlist
\item
  governance
\item
  ICT risk management
\item
  incident handling and reporting
\item
  resilience testing
\item
  recovery capability
\item
  third-party risk management
\item
  concentration risk management
\item
  evidence and reporting
\item
  assurance
\item
  improvement discipline
\end{itemize}

\section{Practical Templates}\label{practical-templates}

\subsection{Suggested Appendices}\label{suggested-appendices}

\begin{itemize}
\tightlist
\item
  Appendix A: ICT risk assessment template
\item
  Appendix B: Major ICT incident classification template
\item
  Appendix C: Incident reporting workflow
\item
  Appendix D: Digital operational resilience testing plan
\item
  Appendix E: Register of information template
\item
  Appendix F: ICT third-party due diligence checklist
\item
  Appendix G: DORA control catalogue
\item
  Appendix H: DORA RACI matrix
\end{itemize}

\section{Customisation Notes}\label{customisation-notes}

This template should be tailored to reflect:

\begin{itemize}
\tightlist
\item
  legal entity structure
\item
  supervisory context
\item
  financial sector segment
\item
  business service model
\item
  technology estate
\item
  cloud and outsourcing profile
\item
  critical and important functions
\item
  operational resilience maturity
\item
  assurance expectations
\end{itemize}

\section{References}\label{references}

Add authoritative references here, for example:

\begin{itemize}
\tightlist
\item
  Regulation (EU) 2022/2554
\item
  DORA implementing and delegated acts
\item
  ESA technical standards and guidance
\item
  competent authority guidance
\item
  internal ICT risk and resilience policies
\item
  outsourcing and third-party risk standards
\item
  enterprise risk framework
\end{itemize}




\end{document}
